\documentclass[12pt, a4paper]{report}

% ─── Packages ───────────────────────────────────────────────────────────────
\usepackage[margin=2.5cm, headheight=25pt, top=2.8cm]{geometry}
\usepackage{graphicx}
\usepackage{xcolor}
\usepackage{titlesec}
\usepackage{titling}
\usepackage{fancyhdr}
\usepackage{enumitem}
\usepackage{listings}
\usepackage{booktabs}
\usepackage{tabularx}
\usepackage{array}
\usepackage{mdframed}
\usepackage{hyperref}
\usepackage{tcolorbox}
\usepackage{tikz}
\usepackage{amsmath}
\usepackage{parskip}
\usepackage[expansion=false]{microtype}
\usepackage{setspace}
\usepackage{caption}
\usepackage{float}

\tcbuselibrary{skins, breakable, listings}
\usetikzlibrary{shapes.geometric, arrows.meta, positioning, fit, backgrounds}

% ─── Colors ─────────────────────────────────────────────────────────────────
\definecolor{farmgreen}{RGB}{34, 120, 60}
\definecolor{soilbrown}{RGB}{101, 67, 33}
\definecolor{skyblue}{RGB}{30, 100, 180}
\definecolor{alertred}{RGB}{180, 30, 30}

% Code Block Specific Colors (Matching the provided light theme)
\definecolor{codebg}{RGB}{244, 244, 240}      % Light beige background
\definecolor{codepink}{RGB}{233, 30, 99}      % Pink for keywords (const, new, null, false)
\definecolor{codepurple}{RGB}{156, 39, 176}   % Purple for strings
\definecolor{codegray}{RGB}{150, 150, 150}    % Gray for line numbers and comments
\definecolor{codetext}{RGB}{0, 0, 0}          % Black for standard code text

% ─── Hyperref Setup ──────────────────────────────────────────────────────────
\hypersetup{
    colorlinks=true,
    linkcolor=farmgreen,
    urlcolor=skyblue,
    citecolor=farmgreen,
    pdfauthor={OOP II Project Team},
    pdftitle={Smart Farm Management System — Project Report},
    pdfsubject={College Project Documentation}
}

% ─── Header & Footer ─────────────────────────────────────────────────────────
\pagestyle{fancy}
\fancyhf{}
\fancyhead[L]{\textcolor{farmgreen}{\textbf{Smart Farm Management System}}}
\fancyhead[R]{\textcolor{gray}{\small OOP II Project Report}}
\fancyfoot[C]{\textcolor{gray}{\thepage}}
\renewcommand{\headrulewidth}{0.4pt}
\renewcommand{\headrule}{\hbox to\headwidth{\color{farmgreen}\leaders\hrule height \headrulewidth\hfill}}

% ─── Section Title Styling ───────────────────────────────────────────────────
\titleformat{\chapter}[block]
  {\normalfont\LARGE\bfseries\color{farmgreen}}
  {\thechapter.}{12pt}{}
  [\vspace{2pt}{\color{farmgreen}\titlerule[1.5pt]}]

\titlespacing*{\chapter}{0pt}{-25pt}{20pt}

\titleformat{\section}
  {\normalfont\large\bfseries\color{soilbrown}}
  {\thesection}{10pt}{}

\titleformat{\subsection}
  {\normalfont\normalsize\bfseries\color{skyblue}}
  {\thesubsection}{8pt}{}

% ─── Code Listing Style ──────────────────────────────────────────────────────
\lstdefinestyle{javaStyle}{
    basicstyle=\color{codetext}\ttfamily\footnotesize,
    identifierstyle=\color{codetext},
    keywordstyle=\color{codepink},
    stringstyle=\color{codepurple},
    commentstyle=\color{codegray}\itshape,
    numberstyle=\tiny\color{codegray},
    numbers=left,
    numbersep=12pt,
    breaklines=true,
    frame=none,
    xleftmargin=24pt,
    xrightmargin=6pt,
    showstringspaces=false,
    tabsize=4,
    language=Java,
    morekeywords={String, Thread, Socket, ServerSocket, Runnable, var, record, SensorHandler, IOException, SQLException, PreparedStatement, ResultSet, Connection, List, ArrayList, LocalDateTime, LocalDate, Date, HashMap, Map, Alert, Crop, Task, Worker, SensorReading, HarvestRecord, Platform, DashboardController, FXCollections, SimpleStringProperty, true, false, null, Stage, Scene, VBox, HBox, Label, Button, GridPane, TextField, Spinner, ComboBox, ToggleButton, Insets, Font, FontWeight, Image, FXMLLoader, Parent, Application, DriverManager, Properties, InputStream, ExecutorService, Executors, BufferedReader, InputStreamReader, GenericDAO},
}

\lstdefinestyle{cppStyle}{
    basicstyle=\color{codetext}\ttfamily\footnotesize,
    identifierstyle=\color{codetext},
    keywordstyle=\color{codepink},
    stringstyle=\color{codepurple},
    commentstyle=\color{codegray}\itshape,
    numberstyle=\tiny\color{codegray},
    numbers=left,
    numbersep=12pt,
    breaklines=true,
    frame=none,
    xleftmargin=24pt,
    xrightmargin=6pt,
    showstringspaces=false,
    tabsize=4,
    language=C++,
    morekeywords={WiFiClient, DHT, String, Serial, true, false, null}
}

\lstdefinestyle{sqlStyle}{
    basicstyle=\color{codetext}\ttfamily\footnotesize,
    identifierstyle=\color{codetext},
    keywordstyle=\color{codepink},
    stringstyle=\color{codepurple},
    commentstyle=\color{codegray}\itshape,
    numberstyle=\tiny\color{codegray},
    numbers=left,
    numbersep=12pt,
    breaklines=true,
    frame=none,
    xleftmargin=24pt,
    xrightmargin=6pt,
    showstringspaces=false,
    tabsize=4,
    language=SQL,
}

% ─── Custom Boxes ─────────────────────────────────────────────────────────────
\tcbset{
  notebox/.style={
    enhanced, breakable,
    colback=farmgreen!8,
    colframe=farmgreen,
    fonttitle=\bfseries,
    coltitle=white,
    attach boxed title to top left={yshift=-2mm, xshift=6mm},
    boxed title style={colback=farmgreen, rounded corners}
  },
  warnbox/.style={
    enhanced,
    colback=alertred!8,
    colframe=alertred,
    fonttitle=\bfseries,
    coltitle=white,
    attach boxed title to top left={yshift=-2mm, xshift=6mm},
    boxed title style={colback=alertred, rounded corners}
  },
  infobox/.style={
    enhanced, breakable,
    colback=skyblue!8,
    colframe=skyblue,
    fonttitle=\bfseries,
    coltitle=white,
    attach boxed title to top left={yshift=-2mm, xshift=6mm},
    boxed title style={colback=skyblue, rounded corners}
  }
}

% Custom tcblisting for exact match to the light-themed code blocks
\newtcblisting[auto counter, number within=chapter]{codeblock}[2]{
    listing only,
    breakable,
    pad at break=2mm,
    bottomsep at break=5pt,
    topsep at break=5pt,
    colback=codebg,
    colframe=codebg,
    coltext=codetext,
    boxrule=0pt,
    arc=0pt,
    outer arc=0pt,
    top=8pt,
    bottom=8pt,
    left=0pt,
    right=8pt,
    listing options={style=#1},
    after={\par\vspace{4pt}\begin{center}Listing \thetcbcounter: #2\end{center}\vspace{16pt}}
}

% ─── Document ─────────────────────────────────────────────────────────────────
\begin{document}

% ══════════════════════════════════════════════════════════════════════════════
%  TITLE PAGE
% ══════════════════════════════════════════════════════════════════════════════
\begin{titlepage}
\begin{tikzpicture}[remember picture, overlay]
    \fill[farmgreen] (current page.north west) rectangle ([yshift=-14cm]current page.north east);
    \fill[soilbrown] (current page.south west) rectangle ([yshift=1.8cm]current page.south east);
    \fill[farmgreen!30, opacity=0.4] ([xshift=12cm, yshift=-4cm]current page.north west) circle (5cm);
\end{tikzpicture}

\vspace*{1.5cm}
\begin{center}
    {\color{white}\fontsize{28}{34}\selectfont\bfseries Smart Farm}\\[4pt]
    {\color{white}\fontsize{22}{28}\selectfont\bfseries Management System}\\[10pt]
    {\color{white!80}\large Project Report}

    \vspace{3.5cm}

    {\color{white}\rule{10cm}{1pt}}\\[18pt]
    {\large\bfseries\color{white} OOP II Module --- Project Report}\\[10pt]
    {\color{white!90} Errors \& Debugging $\cdot$ UML Diagrams $\cdot$ Code Snippets}\\[30pt]
    {\color{white}\rule{10cm}{1pt}}

    \vfill

    \begin{tabular}{rl}
        \textbf{Module:}   & Object-Oriented Programming II \\[4pt]
        \textbf{Hardware:} & ESP32 + DHT11 Sensor + FC-28 + R307 \\[4pt]
        \textbf{Stack:}    & Java, MySQL, JavaFX, Arduino C++ \\[4pt]
        \textbf{Year:}     & 2025 / 2026 \\[12pt]
    \end{tabular}

    \vspace{8pt}
    {\small
    \begin{tabular}{ll}
        Mostafa Mahmoud Aborehab        & 241008538 \\
        Jouiria Mohammed Hassan         & 241008818 \\
        Ahmed Hagag Abdelhamed          & 241002485 \\
        Mohamed Ahmed Mahmoud Abdelbary & 241009209 \\
        Mariam Yasser Dahab             & 241008888 \\
    \end{tabular}}
\end{center}
\end{titlepage}

% ══════════════════════════════════════════════════════════════════════════════
\tableofcontents
\newpage


% ══════════════════════════════════════════════════════════════════════════════
\chapter{Project Information}
% ══════════════════════════════════════════════════════════════════════════════

\section{Group Members}

\begin{table}[H]
\centering
\renewcommand{\arraystretch}{1.5}
\begin{tabularx}{\textwidth}{l l X}
\toprule
\textbf{Name} & \textbf{ID} & \textbf{Role} \\
\midrule
Mostafa Mahmoud Aborehab        & 241008538 & Team Lead, Backend \& IoT \\
Jouiria Mohammed Hassan         & 241008818 & Frontend \& UI \\
Ahmed Hagag Abdelhamed          & 241002485 & Android \& Integration \\
Mohamed Ahmed Mahmoud Abdelbary & 241009209 & Backend \& Database \\
Mariam Yasser Dahab             & 241008888 & UI/UX \& Testing \\
\bottomrule
\end{tabularx}
\caption{Group members and IDs}
\end{table}

\section{Selected Business Domain}

\textbf{Agriculture / Smart Farming.} The system manages a real farm operation --- monitoring field conditions with IoT sensors, tracking crops and plots, assigning workers to tasks, logging harvests, and generating reports. The domain was chosen because it naturally requires every OOP~II concept: multithreading (concurrent sensor data), networking (TCP sockets), database (JDBC/MySQL), GUI (JavaFX), and file I/O (CSV export).

\section{Steps to Run the Project}

\begin{enumerate}[leftmargin=*, label=\textbf{\arabic*.}]
    \item \textbf{Prerequisites:} Install Java~17+, Maven, and MySQL~8.
    \item \textbf{Clone the repository:}
\begin{tcolorbox}[colback=codebg, colframe=codebg, boxrule=0pt, arc=0pt, left=8pt, right=8pt, top=4pt, bottom=4pt]
\texttt{git clone https://github.com/GoBloxy/Agrilliant.git}
\end{tcolorbox}
    \item \textbf{Configure the database} (see Section~\ref{sec:dbsetup} below).
    \item \textbf{Build and run:}
\begin{tcolorbox}[colback=codebg, colframe=codebg, boxrule=0pt, arc=0pt, left=8pt, right=8pt, top=4pt, bottom=4pt]
\texttt{cd Agrilliant}\\
\texttt{mvn clean javafx:run}
\end{tcolorbox}
    \item The application opens the sign-in screen. Log in with your credentials or use fingerprint scan for worker access.
    \item The TCP sensor server starts automatically on port \textbf{8080} (see Section~\ref{sec:socket}).
\end{enumerate}

\section{Database Setup}\label{sec:dbsetup}

The project uses a \textbf{MySQL~8} database. Connection details are stored in \texttt{src/main/resources/db.properties}:

\begin{tcolorbox}[colback=codebg, colframe=codebg, boxrule=0pt, arc=0pt, left=8pt, right=8pt, top=6pt, bottom=6pt]
\ttfamily\small
db.url=jdbc:mysql://\textit{<host>}:3306/agrilliant\\
db.user=\textit{<username>}\\
db.password=\textit{<password>}
\end{tcolorbox}

\textbf{Setup steps:}
\begin{enumerate}[noitemsep]
    \item Create a MySQL database named \texttt{agrilliant}.
    \item The application uses a cloud-hosted MySQL instance by default. To use a local database, update \texttt{db.properties} with your local MySQL credentials.
    \item Tables are created automatically by the application on first run, or you can run the schema script from the project documentation.
\end{enumerate}

\section{Socket Port Number}\label{sec:socket}

The TCP sensor server listens on port \textbf{8080}. ESP32 devices connect to this port to send live sensor data.

\begin{table}[H]
\centering
\renewcommand{\arraystretch}{1.5}
\begin{tabularx}{\textwidth}{l l X}
\toprule
\textbf{Service} & \textbf{Port} & \textbf{Description} \\
\midrule
TCP Sensor Server & 8080 & Receives live data from ESP32 devices \\
MQTT Broker       & 1883 & Distributes sensor data to all clients (HiveMQ public broker) \\
MySQL Database    & 3306 & Persistent storage for all farm data \\
\bottomrule
\end{tabularx}
\caption{Network ports used by the system}
\end{table}

\begin{tcolorbox}[warnbox, title={Important}]
Port 8080 must be available on the machine running the server. If another application is using this port, the TCP server will print a warning but the application will still function --- MQTT will deliver sensor data independently.
\end{tcolorbox}


% ══════════════════════════════════════════════════════════════════════════════
\chapter{Errors \& Debugging}
% ══════════════════════════════════════════════════════════════════════════════

% TODO: For EACH error include:
%   - Error description
%   - Screenshot of the error
%   - Cause of the error
%   - Solution
%   - Related code snippet

\textit{This section will document errors encountered during development, their root causes, and how they were resolved.}


% ══════════════════════════════════════════════════════════════════════════════
\chapter{UML Diagrams}
% ══════════════════════════════════════════════════════════════════════════════

% TODO: Class Diagram (REQUIRED)

\textit{This section will contain the UML class diagram for the system.}


% ══════════════════════════════════════════════════════════════════════════════
\chapter{Code Snippets}
% ══════════════════════════════════════════════════════════════════════════════

This chapter presents representative code snippets from the Smart Farm system, each demonstrating a core OOP~II concept.


% ──────────────────────────────────────────────────────────────────────────────
\section{Class, Constructor \& Methods}
% ──────────────────────────────────────────────────────────────────────────────

The \texttt{Alert} class models a farm alert triggered by sensor readings. It demonstrates:

\begin{itemize}[noitemsep]
    \item An \textbf{inner enum} (\texttt{Severity}) to restrict field values
    \item \textbf{Multiple constructors} (overloading) for different creation contexts
    \item \textbf{Instance methods} that encapsulate behaviour (\texttt{resolve()}, \texttt{isCritical()})
    \item \textbf{Encapsulation} with private fields and public getters/setters
\end{itemize}

\begin{codeblock}{javaStyle}{Alert --- class with enum, constructors, and methods}
public class Alert {
    public enum Severity { INFO, WARNING, CRITICAL }

    private int           alertId;
    private String        alertType;
    private Severity      severity;
    private String        message;
    private boolean       resolved;
    private LocalDateTime timestamp;
    private int           plotId;

    // Full constructor (used when loading from database)
    public Alert(int alertId, String alertType, Severity severity,
                 String message, boolean resolved,
                 LocalDateTime timestamp, int plotId) {
        this.alertId   = alertId;
        this.alertType = alertType;
        this.severity  = severity;
        this.message   = message;
        this.resolved  = resolved;
        this.timestamp = timestamp;
        this.plotId    = plotId;
    }

    // Convenience constructor (used when creating a new alert)
    public Alert(String alertType, Severity severity,
                 String message, int plotId) {
        this.alertId   = -1;
        this.alertType = alertType;
        this.severity  = severity;
        this.message   = message;
        this.resolved  = false;
        this.timestamp = LocalDateTime.now();
        this.plotId    = plotId;
    }

    // Behaviour methods
    public void    resolve()    { this.resolved = true; }
    public boolean isCritical() { return severity == Severity.CRITICAL; }
    public boolean isResolved() { return resolved; }

    // Getters
    public int           getAlertId()   { return alertId; }
    public String        getAlertType() { return alertType; }
    public Severity      getSeverity()  { return severity; }
    public String        getMessage()   { return message; }
    public LocalDateTime getTimestamp()  { return timestamp; }
    public int           getPlotId()    { return plotId; }

    @Override
    public String toString() {
        return "[" + severity + "] " + alertType
             + " -- " + message + " (plot=" + plotId + ")";
    }
}
\end{codeblock}

\begin{tcolorbox}[notebox, title={Key OOP Concepts}]
\begin{itemize}[noitemsep]
    \item \textbf{Encapsulation} --- all fields are \texttt{private}; access is through methods.
    \item \textbf{Constructor overloading} --- two constructors serve different use cases (DB load vs.\ new alert).
    \item \textbf{Enum} --- \texttt{Severity} restricts the value to exactly three options, enforced at compile time.
\end{itemize}
\end{tcolorbox}


% ──────────────────────────────────────────────────────────────────────────────
\section{JavaFX Stage (Programmatic --- No FXML)}
% ──────────────────────────────────────────────────────────────────────────────

The application entry point creates the \texttt{Stage}, configures the window, and launches background services. This is pure Java --- no FXML layout file is used for the Stage itself.

\begin{codeblock}{javaStyle}{Main --- JavaFX Application with programmatic Stage setup}
public class Main extends Application {

    @Override
    public void start(Stage primaryStage) throws Exception {
        // Build the root layout in Java code
        Parent root = FXMLLoader.load(
            getClass().getResource("/fxml/signin.fxml"));

        // Create Scene and apply stylesheet
        Scene scene = new Scene(root, 1000, 650);
        scene.getStylesheets().add(
            getClass().getResource("/css/farm-theme.css")
                      .toExternalForm());

        // Configure the Stage (window)
        primaryStage.setTitle(
            "Agrilliant -- Smart Farm Management System");
        primaryStage.getIcons().add(new Image(
            getClass().getResourceAsStream("/images/logo-dark.png")));
        primaryStage.setScene(scene);
        primaryStage.setMinWidth(1300);
        primaryStage.setMinHeight(820);
        primaryStage.setMaximized(true);
        primaryStage.setOnCloseRequest(e -> System.exit(0));
        primaryStage.show();
    }

    public static void main(String[] args) {
        launch(args);
    }
}
\end{codeblock}

The \texttt{SettingsPage} is a screen built \textbf{entirely without FXML} --- all UI controls are created and laid out in pure Java:

\begin{codeblock}{javaStyle}{SettingsPage --- programmatic JavaFX UI (no FXML)}
public class SettingsPage extends VBox {

    private final TextField txtServerIp   = new TextField("192.168.8.141");
    private final TextField txtServerPort = new TextField("8080");
    private final Spinner<Integer> spnInterval = new Spinner<>(5, 300, 60, 5);
    private final ComboBox<String> cmbUnit     = new ComboBox<>();

    public SettingsPage() {
        setSpacing(24);
        setPadding(new Insets(32, 40, 32, 40));

        // Header
        Label title = new Label("Settings");
        title.setFont(Font.font("System", FontWeight.BOLD, 24));

        // Server section
        GridPane serverGrid = new GridPane();
        serverGrid.setHgap(16);
        serverGrid.setVgap(12);
        serverGrid.add(new Label("Server IP"), 0, 0);
        serverGrid.add(txtServerIp, 1, 0);
        serverGrid.add(new Label("Port"), 2, 0);
        serverGrid.add(txtServerPort, 3, 0);

        // Temperature unit selector
        cmbUnit.getItems().addAll("Celsius", "Fahrenheit");
        cmbUnit.setValue("Celsius");

        // Save button with event handler
        Button btnSave = new Button("Save Changes");
        btnSave.setOnAction(e -> {
            Alert alert = new Alert(Alert.AlertType.INFORMATION,
                "Settings saved.", ButtonType.OK);
            alert.showAndWait();
        });

        // Assemble the layout
        getChildren().addAll(title, serverGrid, cmbUnit, btnSave);
    }
}
\end{codeblock}

\begin{tcolorbox}[infobox, title={Why Two Approaches?}]
JavaFX supports two UI-building strategies:
\begin{itemize}[noitemsep]
    \item \textbf{FXML} --- XML layout files (used for complex screens like the dashboard)
    \item \textbf{Programmatic} --- pure Java code (used for simpler pages like Settings)
\end{itemize}
The \texttt{Stage} and \texttt{Scene} are always created in Java code. FXML only defines the \emph{content} inside the scene.
\end{tcolorbox}


% ──────────────────────────────────────────────────────────────────────────────
\section{JDBC Database Programming}
% ──────────────────────────────────────────────────────────────────────────────

\subsection{Database Connection (Singleton)}

The \texttt{DBConnection} class provides a single shared \texttt{Connection} instance using the Singleton pattern. Credentials are loaded from a properties file.

\begin{codeblock}{javaStyle}{DBConnection --- Singleton JDBC connection}
public class DBConnection {
    private static Connection connection;
    private static String url, user, password;

    private DBConnection() {} // prevent instantiation

    static {
        // Load credentials from db.properties
        try (InputStream input = DBConnection.class
                .getClassLoader()
                .getResourceAsStream("db.properties")) {
            Properties props = new Properties();
            props.load(input);
            url      = props.getProperty("db.url");
            user     = props.getProperty("db.user");
            password = props.getProperty("db.password");
        } catch (IOException e) {
            System.err.println("Failed to load db.properties");
        }
    }

    public static Connection getInstance() {
        try {
            if (connection == null || connection.isClosed()) {
                connection = DriverManager.getConnection(
                    url, user, password);
            }
        } catch (SQLException e) {
            System.err.println("DB connection failed: "
                + e.getMessage());
        }
        return connection;
    }
}
\end{codeblock}

\subsection{CRUD with PreparedStatement}

The \texttt{WorkerDAO} demonstrates JDBC CRUD operations using \texttt{PreparedStatement} to prevent SQL injection:

\begin{codeblock}{javaStyle}{WorkerDAO --- INSERT and SELECT with PreparedStatement}
public class WorkerDAO implements GenericDAO<Worker> {
    private final Connection conn = DBConnection.getInstance();

    @Override
    public void save(Worker item) throws SQLException {
        String sql = "INSERT INTO worker (full_name, phone, email, "
                   + "job_title, on_duty) VALUES (?, ?, ?, ?, ?)";
        try (PreparedStatement ps = conn.prepareStatement(sql,
                Statement.RETURN_GENERATED_KEYS)) {
            ps.setString (1, item.getFullName());
            ps.setString (2, item.getPhone());
            ps.setString (3, item.getEmail());
            ps.setString (4, item.getJobTitle());
            ps.setBoolean(5, item.isOnDuty());
            ps.executeUpdate();

            // Retrieve auto-generated primary key
            try (ResultSet keys = ps.getGeneratedKeys()) {
                if (keys.next())
                    item.setWorkerId(keys.getInt(1));
            }
        }
    }

    @Override
    public List<Worker> getAll() throws SQLException {
        String sql = "SELECT * FROM worker ORDER BY full_name";
        List<Worker> workers = new ArrayList<>();
        try (PreparedStatement ps = conn.prepareStatement(sql);
             ResultSet rs = ps.executeQuery()) {
            while (rs.next()) {
                workers.add(new Worker(
                    rs.getInt("worker_id"),
                    rs.getString("full_name"),
                    rs.getString("phone"),
                    rs.getString("email"),
                    rs.getString("job_title"),
                    rs.getBoolean("on_duty")
                ));
            }
        }
        return workers;
    }

    @Override
    public void delete(int id) throws SQLException {
        String sql = "DELETE FROM worker WHERE worker_id = ?";
        try (PreparedStatement ps = conn.prepareStatement(sql)) {
            ps.setInt(1, id);
            ps.executeUpdate();
        }
    }
}
\end{codeblock}

\begin{tcolorbox}[notebox, title={JDBC Best Practices Shown}]
\begin{itemize}[noitemsep]
    \item \textbf{PreparedStatement} --- prevents SQL injection by using \texttt{?}~placeholders instead of string concatenation.
    \item \textbf{try-with-resources} --- automatically closes \texttt{PreparedStatement} and \texttt{ResultSet}, even if an exception is thrown.
    \item \textbf{Generated keys} --- retrieves the auto-incremented ID after an \texttt{INSERT}.
\end{itemize}
\end{tcolorbox}


% ──────────────────────────────────────────────────────────────────────────────
\section{Multithreading}
% ──────────────────────────────────────────────────────────────────────────────

The system uses multithreading in two ways: a \textbf{daemon thread} for the TCP server and an \textbf{ExecutorService thread pool} that spawns a new thread for each connected sensor device.

\subsection{Starting the Server on a Daemon Thread}

The TCP server blocks forever waiting for connections, so it must run on a background thread. A \textbf{daemon thread} is used so it automatically stops when the application closes.

\begin{codeblock}{javaStyle}{Main --- launching the server on a daemon thread}
// Inside Main.start() -- after showing the JavaFX stage:

// Start TCP sensor server on a daemon thread
Thread serverThread = new Thread(() -> new FarmServer().start());
serverThread.setDaemon(true);       // dies when app closes
serverThread.setName("FarmServer-TCP");
serverThread.start();
\end{codeblock}

\subsection{Thread Pool for Multiple Devices}

The \texttt{FarmServer} uses a \texttt{CachedThreadPool} to handle multiple ESP32 devices simultaneously. Each device gets its own \texttt{SensorHandler} thread.

\begin{codeblock}{javaStyle}{FarmServer --- thread pool accepts multiple device connections}
public class FarmServer {
    private static final int PORT = 8080;

    public void start() {
        ExecutorService threadPool = Executors.newCachedThreadPool();
        try (ServerSocket serverSocket = new ServerSocket(PORT)) {
            System.out.println("Farm Server started on port " + PORT);

            while (true) {
                Socket deviceSocket = serverSocket.accept();
                System.out.println("Device connected: "
                    + deviceSocket.getInetAddress());
                // Each device runs on its own thread
                threadPool.submit(new SensorHandler(deviceSocket));
            }
        } catch (IOException e) {
            System.err.println("Server error: " + e.getMessage());
        }
    }
}
\end{codeblock}

\begin{tcolorbox}[infobox, title={Multithreading Concepts Demonstrated}]
\begin{itemize}[noitemsep]
    \item \textbf{Daemon thread} --- background thread that does not prevent JVM shutdown.
    \item \textbf{ExecutorService} --- manages a pool of threads; avoids creating threads manually for each connection.
    \item \textbf{Runnable} --- \texttt{SensorHandler} implements \texttt{Runnable}, allowing it to be submitted to the thread pool.
    \item \textbf{Concurrency} --- multiple ESP32 sensors can connect and send data simultaneously without blocking each other.
\end{itemize}
\end{tcolorbox}


% ──────────────────────────────────────────────────────────────────────────────
\section{Socket Programming}
% ──────────────────────────────────────────────────────────────────────────────

The system uses TCP sockets to receive live data from ESP32 devices. The \textbf{server side} uses \texttt{ServerSocket} to listen for connections, and the \textbf{handler} uses \texttt{Socket.getInputStream()} to read data line by line.

\begin{codeblock}{javaStyle}{SensorHandler --- reads data from a TCP socket}
public class SensorHandler implements Runnable {
    private final Socket socket;
    private final SensorService sensorService = new SensorService();

    public SensorHandler(Socket socket) {
        this.socket = socket;
    }

    @Override
    public void run() {
        try (BufferedReader reader = new BufferedReader(
                new InputStreamReader(socket.getInputStream()))) {
            String line;
            while ((line = reader.readLine()) != null) {
                SensorReading reading = parseLine(line);
                if (reading != null) {
                    sensorService.processReading(reading);
                }
            }
        } catch (IOException e) {
            System.out.println("Device disconnected: "
                + socket.getInetAddress());
        }
    }

    // Parses "DEVICE:plot1_sensor,TEMP:27.50,HUM:63.20"
    private SensorReading parseLine(String raw) {
        try {
            String[] parts = raw.split(",");
            String device = parts[0].split(":")[1];
            float  temp   = Float.parseFloat(parts[1].split(":")[1]);
            float  hum    = Float.parseFloat(parts[2].split(":")[1]);
            return new SensorReading(device, temp, hum,
                                     LocalDateTime.now());
        } catch (Exception e) {
            System.err.println("Bad data format: " + raw);
            return null;
        }
    }
}
\end{codeblock}

The data flow is:

\begin{tcolorbox}[colback=codebg, colframe=gray, fontupper=\ttfamily\small\color{codetext}]
\begin{verbatim}
ESP32 (WiFiClient)
  |--- TCP connect ---> ServerSocket.accept()
  |--- "DEVICE:plot1_sensor,TEMP:27.5,HUM:63.2\n" --->
  |                     BufferedReader.readLine()
  |                         |
  |                     parseLine() -> SensorReading
  |                         |
  |                     sensorService.processReading()
  |                         |
  |                     Save to DB + Check alerts + Update UI
\end{verbatim}
\end{tcolorbox}

\begin{tcolorbox}[notebox, title={Socket Programming Concepts}]
\begin{itemize}[noitemsep]
    \item \textbf{ServerSocket} --- listens on a port and accepts incoming connections.
    \item \textbf{Socket} --- represents one end of a two-way TCP connection.
    \item \textbf{InputStream wrapping} --- \texttt{Socket} $\rightarrow$ \texttt{InputStreamReader} $\rightarrow$ \texttt{BufferedReader} for line-by-line reading.
    \item \textbf{Protocol design} --- a simple text-based protocol (\texttt{KEY:VALUE,KEY:VALUE}) parsed with \texttt{split()}.
\end{itemize}
\end{tcolorbox}


% ──────────────────────────────────────────────────────────────────────────────
\section{Generics}
% ──────────────────────────────────────────────────────────────────────────────

The system uses a \textbf{generic DAO interface} to define standard CRUD operations for \emph{any} model class. Each concrete DAO implements this interface with a specific type.

\subsection{Generic Interface Definition}

\begin{codeblock}{javaStyle}{GenericDAO<T> --- generic interface for all data access objects}
public interface GenericDAO<T> {
    void    save(T item)    throws SQLException;
    T       getById(int id) throws SQLException;
    List<T> getAll()        throws SQLException;
    void    update(T item)  throws SQLException;
    void    delete(int id)  throws SQLException;
}
\end{codeblock}

\subsection{Concrete Implementation}

Each model class gets its own DAO that implements \texttt{GenericDAO} with the appropriate type parameter:

\begin{codeblock}{javaStyle}{TaskDAO --- implements GenericDAO<Task>}
public class TaskDAO implements GenericDAO<Task> {
    private final Connection conn = DBConnection.getInstance();

    @Override
    public void save(Task item) throws SQLException {
        String sql = "INSERT INTO tasks (description, status, "
                   + "due_date, plot_id) VALUES (?, ?, ?, ?)";
        try (PreparedStatement ps = conn.prepareStatement(sql,
                Statement.RETURN_GENERATED_KEYS)) {
            ps.setString(1, item.getDescription());
            ps.setString(2, item.getStatus().name());
            ps.setDate  (3, Date.valueOf(item.getDueDate()));
            ps.setInt   (4, item.getPlotId());
            ps.executeUpdate();

            try (ResultSet keys = ps.getGeneratedKeys()) {
                if (keys.next())
                    item.setTaskId(keys.getInt(1));
            }
        }
    }

    @Override
    public List<Task> getAll() throws SQLException {
        List<Task> tasks = new ArrayList<>();
        String sql = "SELECT * FROM tasks ORDER BY due_date";
        try (PreparedStatement ps = conn.prepareStatement(sql);
             ResultSet rs = ps.executeQuery()) {
            while (rs.next()) {
                tasks.add(new Task(
                    rs.getInt("task_id"),
                    rs.getString("description"),
                    Task.Status.valueOf(rs.getString("status")),
                    rs.getDate("due_date").toLocalDate(),
                    rs.getInt("plot_id")
                ));
            }
        }
        return tasks;
    }

    @Override
    public void delete(int id) throws SQLException {
        String sql = "DELETE FROM tasks WHERE task_id = ?";
        try (PreparedStatement ps = conn.prepareStatement(sql)) {
            ps.setInt(1, id);
            ps.executeUpdate();
        }
    }

    // getById() and update() follow the same pattern...
}
\end{codeblock}

All DAO classes in the project implement this single generic interface:

\begin{table}[H]
\centering
\renewcommand{\arraystretch}{1.5}
\begin{tabularx}{\textwidth}{>{\ttfamily}l >{\ttfamily}l X}
\toprule
\textbf{DAO Class} & \textbf{Implements} & \textbf{Model} \\
\midrule
TaskDAO    & GenericDAO<Task>          & Farm tasks assigned to workers \\
WorkerDAO  & GenericDAO<Worker>        & Farm workers \\
AlertDAO   & GenericDAO<Alert>         & Sensor-triggered alerts \\
CropDAO    & GenericDAO<Crop>          & Planted crops \\
PlotDAO    & GenericDAO<Plot>          & Farm plots of land \\
DeviceDAO  & GenericDAO<Device>        & IoT sensor devices \\
\bottomrule
\end{tabularx}
\caption{All DAO classes implement the same generic interface}
\end{table}

\begin{tcolorbox}[notebox, title={Generics Concepts Demonstrated}]
\begin{itemize}[noitemsep]
    \item \textbf{Type parameter \texttt{<T>}} --- the interface is written once and works for any model type.
    \item \textbf{Type safety} --- \texttt{GenericDAO<Task>} guarantees that \texttt{save()} only accepts a \texttt{Task}, not a \texttt{Worker} or \texttt{Alert}. The compiler catches type mismatches.
    \item \textbf{Code reuse} --- the contract (five CRUD methods) is defined once; six classes implement it without duplicating the interface.
    \item \textbf{Bounded usage} --- \texttt{List<T>} in the return type of \texttt{getAll()} is itself a generic, showing generics composed inside generics.
\end{itemize}
\end{tcolorbox}


\end{document}
